\chapter{Whose words these are}

These languages belong to the people who speak them. None of this work exists without them. Everything else in this research is measured against the tables below, and the tables are their words, written down.

Each entry opens with who spoke, because that is whose language it is. A linguist wrote the paper and a person read the paper into a table, and neither of those is whose language it is. Where a paper cites a published dictionary and never says who spoke, the entry says so, and the linguist does not go in the speaker column.

Conditions the speakers set are recorded with them below and hold wherever this corpus is used.

\subsection{Lushootseed}

\begin{itemize}
\item \textbf{Martha Lamont, Northern dialect}
\item \textbf{Annie Jack, Southern dialect}
\end{itemize}

\cite{src:Mellesmoen-and-Kye-2026}, 467 rows.

Both recorded by Leon Metcalf in the 1950s. The only paper here that labels every form by dialect, and the border test is scored against it.

\subsection{Lushootseed}

\begin{itemize}
\item \textbf{Susie Sampson Peter, Upper Skagit}
\item \textbf{Martha LaMont, Tulalip-Skagit}
\end{itemize}

\cite{src:Hilbert-1983}, 60 rows.

Vi taqʷšəblu Hilbert wrote the paper; the twenty-one examples were said by her aunt Susie Sampson Peter and by Martha LaMont, recorded by Leon Metcalf between 1950 and 1958 and by Thom Hess in 1963. Hilbert is the author of the paper and not its speaker.

\subsection{St'át'imcets}

\begin{itemize}
\item \textbf{K̓weswapáw̓ (Linda Redan), Qayqáyten}
\item \textbf{Sam Mitchell, in van Eijk and Williams 1981}
\end{itemize}

\cite{src:Matthewson-and-Redan-2026}, 104 rows.

K̓weswapáw̓ told the story over Zoom on 31 October 2025, three minutes twenty-eight seconds, and the audio and video are held by her. Sam Mitchell is the speaker of the earlier text the paper cites.

\subsection{St'át'imcets}

\begin{itemize}
\item \textbf{Qwa7yán'ak (Carl Alexander), Nxwísten}
\end{itemize}

\cite{src:Alexander-and-Davis-2026}, 691 rows.

Recorded at Nxwísten on 7 July 2025, just over half an hour.

\subsection{Nuxalk}

\begin{itemize}
\item \textbf{Dr. Margaret Siwallace}
\end{itemize}

\cite{src:Nater-2015}, 285 rows.

Recorded about 1975, published 2015.

\subsection{nɬeʔkepmxcín}

\begin{itemize}
\item \textbf{wlwlmelst (Maurice Michell), Southern yutémkt dialect}
\end{itemize}

\cite{src:LaFontaine-and-Janzen-2024}, 226 rows.

He shares his four stories freely for people connecting with the language. They came from his mother nxwelinek and his grandmother ʔústko.

\subsection{nɬeʔkepmxcín}

\begin{itemize}
\item \textbf{Kʷəɬtəzétkʷu (Bernice Garcia), c̓əɬétkʷu (Coldwater)}
\end{itemize}

\cite{src:Garcia-Hannon-and-Stacey-2024}, 371 rows.

She asks it be acknowledged she is a Kamloops Indian Residential School speaker re-learning her language.

\subsection{Mainland Comox (ayajuthem)}

\begin{itemize}
\item \textbf{Mary George, Sliammon}
\item \textbf{Noel George Harry}
\item \textbf{Tommy Paul}
\end{itemize}

\cite{src:John-Hamilton-Davis-2021}, 1661 rows.

Recorded 1969 to 1980. The two checks of its extraction repair it differently, and their results for it are not directly comparable.

\subsection{nɬeʔkepmxcín}

\begin{itemize}
\item \textbf{Bev Phillips, Lytton First Nation (ƛ̓q̓əmcín)}
\end{itemize}

\cite{src:Hall-and-Phillips-2025}, 406 rows.

Her own recorded reading of the story is held with the corpus, and it checks the extraction as well as sourcing it. The two checks of its extraction close spaces differently.

\subsection{Nsyilxcən}

\begin{itemize}
\item \textbf{George Lezard, Penticton Indian Reserve}
\item \textbf{Nellie Guitterez, Upper Nicola Indian Band}
\item \textbf{Kiláwnaʔ (Andrew McGinnis), Penticton Indian Reserve}
\end{itemize}

\cite{src:Lyon-2015}, 1417 rows.

George Lezard recorded 1966 by Randy Bouchard, transcribed by Larry Pierre 1970, updated by permission of Arnie Baptiste, his son. Nellie Guitterez recorded 1978 or 1979 by Yvonne Hébert, reprinted by permission of Lynne Jorgesen, her great-granddaughter.

\subsection{Nsyilxcən}

\begin{itemize}
\item \textbf{Lottie Lindley, Upper Nicola}
\end{itemize}

\cite{src:Lindley-and-Lyon-2013}, 653 rows.

\subsection{Lushootseed}

\begin{itemize}
\item The paper names no speaker. Its forms are cited from a published source.
\end{itemize}

\cite{src:Hilbert-and-Hess-1975}, 74 rows.

A 1975 typescript, scanned, and its text is OCR of the scan. The paper names no speaker for its examples. The OCR carries none of the orthography: zero schwas, zero raised w, zero barred l, zero wedges, against 169, 77, 41 and 33 in the hand extraction. It writes taqWsablu for taqʷšəblu and slahal for sləhal. All 200 disagreements in which the extraction holds a form the hand extraction does not come from that, and they are faults of the source text and not of the transcriber. Until a page text is drafted for it, as for the two Lyon papers, a check against it measures nothing.

\subsection{nɬeʔkepmxcín and Secwepemctsín}

\begin{itemize}
\item \textbf{Charley Alexis Mayoos}
\item \textbf{William Celestin}
\end{itemize}

\cite{src:Robertson-2012}, 195 rows.

Their texts are written in Chinuk pipa. Texts 3 to 6 are Chinook Jargon, which is a pidgin and is not Salish, and the speaker column keeps those out.

\subsection{eighteen Central Salish languages}

\begin{itemize}
\item The paper names no speaker. Its forms are cited from a published source.
\end{itemize}

\cite{src:Wolfe-2025}, 772 rows.

Every form is cited from a published dictionary of one of eighteen languages. there is nobody this corpus is of. The speaker column carries the language instead, and the gold standard corpus for this paper is kept per language.

\subsection{Nuxalk}

\begin{itemize}
\item The paper names no speaker. Its forms are cited from a published source.
\end{itemize}

\cite{src:Nater-2024}, 671 rows.

Nater's own records, from his 1990 dictionary and 1984 grammar. No speaker is named. Six entries and two tables are Heiltsuk, Oowekyala, Kwak̓wala and Haisla, which are North Wakashan and not Salish at all.

\subsection{Nsyilxcən}

\begin{itemize}
\item \textbf{ɬk̓mxnalqs (Delphine Derrickson-Armstrong), stq̓aʔtkʷɬniw̓t}
\item \textbf{c̓əskʕáknaʔ (Dave Michele), stq̓aʔtkʷɬniw̓t}
\end{itemize}

\cite{src:Lyon-2025}, 537 rows.

Elicited from both speakers. Most cells of its two tables are starred, which is a form the linguist built and the speakers rejected, and those are held out.

\subsection{Twana}

\begin{itemize}
\item The paper names no speaker. Its forms are cited from a published source.
\end{itemize}

\cite{src:Kim-2017}, 358 rows.

Every Twana form is Drachman's, out of a 1969 dissertation the paper calls the only reliable reference in existence for this. No speaker is named. Footnote 9's four Moses-Columbian forms are the only place its extraction repeats a combining mark, and those are corrected from the page.

\subsection{Nuxalk}

\begin{itemize}
\item The paper names no speaker. Its forms are cited from a published source.
\end{itemize}

\cite{src:Nater-2013}, 5767 rows.

1407 numbered entries out of Nater's own 1990 dictionary. No speaker is named.

\subsection{nɬeʔkepmxcín}

\begin{itemize}
\item \textbf{Bev Phillips}
\item \textbf{c̓úʔsinek (Marty Aspinall)}
\item \textbf{kʷaɬtèzetkʷ (Bernice Garcia)}
\end{itemize}

\cite{src:Hall-Luntzlara-Mellesmoen-and-Reid-2026}, 494 rows.

The forms are cited from Thompson and Thompson's grammar and dictionary. These three are the speakers the paper thanks, two examples are Bev Phillips reading her own story, and kʷaɬtèzetkʷ introduces herself in the acknowledgement.

\subsection{St'át'imcets}

\begin{itemize}
\item \textbf{Qwa7yán'ak (Carl Alexander), Nxwísten}
\end{itemize}

\cite{src:Davis-and-Mellesmoen-2023}, 441 rows.

Its data has three sources: van Eijk's dictionary, Davis et al. in preparation, and elicitation with Carl Alexander. It labels forms (U) and (L) for Upper and Lower St'át'imcets, the second external dialect label in the archive.

\subsection{Tillamook}

\begin{itemize}
\item The paper names no speaker. Its forms are cited from a published source.
\end{itemize}

\cite{src:Hamp-1967}, 154 rows.

Eric P. Hamp, Another Look at Tillamook Phonology, ICSNL 2. It cites no words. its language content is four consonant inventories set side by side and two feature matrices. No speaker is named, and the paper works throughout from Thompson and Thompson, Reichard, Kinkade, Drachman and Edel. What it does record is that a Tillamook speaker was living in 1967 and it does not say who: the urgency of the matter is put as being that there is yet a surviving speaker available for possible re-check, and page 2 grants one sense of its claim only on the assumption that a last remaining speaker is typical of a community. Both sentences are in the table. The Twana chart is used, in the page's own words, without his permission, meaning Drachman's. The dot under the uvular fricatives prints solid in some cells of these charts and as an open ring in others, and the two positions swap between the Tillamook and Twana charts. It is one mark and the variation is the typewriter.

\subsection{nɬeʔkepmxcín}

\begin{itemize}
\item \textbf{Bev Phillips, Lytton (ƛ̓q̓əmcín) dialect}
\end{itemize}

\cite{src:Givens-and-Hall-2023}, 178 rows.

Katherine Givens and Brent Hall, The Moon and the Birchbark Canoe (ɬ máʕxetn pe ɬ qʷɬinéwɬ), ICSNL 58. Seven pages, read one at a time off page renders. Page 1: the story was recounted in Nɬeʔkepmxcín by Bev Phillips, a native speaker of the Lytton dialect, who also helped with the translation, and Givens and Hall transcribed and glossed it. She is quoted on the page choosing the story and saying creation stories are not just stories to us, and the paper's first footnote thanks her for entrusting it to them. Of the three recordings, only this one has a file named for the transcribers and not the speaker, and read from that name alone it suggests a second speaker. Footnote 4 defines (VG), a volunteered gloss, as a translated sentence BP offered. The parenthesis at the right margin is the only thing separating her English from the authors', and examples 12, 13 and 14 lack it. Its section 4 repeats every sentence of section 2 in morpheme-broken form, which gave a second independent reading of every word and turned up three places where the two tiers disagree. Footnotes 2, 6, 7 and 8 each declare a mark or a parsing the authors reached by ear or by inference, which is more than any other paper here states about its own readings.

\subsection{Snohomish}

\begin{itemize}
\item The paper names no speaker. Its forms are cited from a published source.
\end{itemize}

\cite{src:Hess-1967}, 248 rows.

Thom Hess, The Morph /-(ə)b/ in Snohomish, ICSNL 2. The hand extraction is checked against the paper. No speaker is named. The forms are Hess's Snohomish data. Marked characters, page 3, by eye and inside one table: xáyəb 'laugh' and xʷúyub 'sell' print a bare x where ƛ̓áɬəb 'salty' prints a barred x body with a glottalization mark, t̓ádəb 'bitter' prints the same t as the English word taste on the first line of that page with a glottalization mark added, and d̓áƛ̓əb 'cloud' carries a hook at the top of its d and a bar across the foot, one composite, where the d of t̓ádəb one column away is bare. Below sixteen times magnification either half of that composite drops out. Glyphs in the printed text unclear. Hilbert and Hess (1975) set the same character in dəxʷgʷəƛ̓əlads, and Nater's Coast Salish *ƛ̓aɬ 'bitter, salt' matches ƛ̓áɬəb segment for segment.

\subsection{Twana}

\begin{itemize}
\item The paper names no speaker. Its forms are cited from a published source.
\end{itemize}

\cite{src:Elmendorf-1967}, 276 rows.

William W. Elmendorf, Word Tabu and Change Rates, ICSNL 2. Thirteen languages are cited in it and Twana is its subject. That is the language named here; every row carries its own language in the speaker column. No speaker is named anywhere: the forms come from Boas and Haeberlin 1927, Krueger 1967, McIlwraith 1948, Walters 1938 and Ray 1932, from Warren Snyder's Suquamish list and Wayne Suttles' Squamish field notes, and from Elmendorf's own field data. This typescript writes its uvular as x under a dot where Hilbert and Hess (1975) write it under a caron, and Hess (1967) agrees with it, which is two papers of one conference against one of a later year. One mark is unresolved and marked so in the table: a short raised stroke over the s of the Columbia ska'u on page 7, which is not the wedge the Upper Chehalis sča'u carries two words earlier.

\subsection{Columbian}

\begin{itemize}
\item The paper names no speaker. Its forms are cited from a published source.
\end{itemize}

\cite{src:Kinkade-1967}, 316 rows.

M. Dale Kinkade, Deictics in Columbian: A Work Paper, ICSNL 2. Twelve pages, read one at a time off page renders. Columbian is what it is about; Kalispel comes from Vogt 1940, Coeur d'Alene from Reichard 1938, and three Colville forms from an unnamed speaker who also knew Columbian. Every row carries its own language in the speaker column. No speaker is named anywhere in the paper: it says my Cm informants and says no more than that. The forms in this table were said by people it does not identify. Its typewriter has three raised marks and the table's symbol notes turn on telling them apart, a comma with a thick head and a curling tail, a V wedge, and a straight acute. čén̓ on page 10 carries all three in one word and is the control. The wedge appears only in the Kalispel and Coeur d'Alene forms, ten times for ten, and never in the Columbian ones, whose seven caron readings are assigned and are all x̌ or the č of čiílx. Whether that split is about the language or about page order is not decidable here, because the wedge first appears on page 8 and every Columbian form was typed on pages 1 to 7. Elmendorf (1967), from the same conference, distinguishes a wedge from a short raised stroke in its own table. the distinction is one these typescripts can carry.

25 tables, 16822 rows read by hand.

